\documentclass[10pt,a4paper]{amsart}
\usepackage[margin=19mm]{geometry}
\usepackage{amsmath,amssymb,mathtools}
\usepackage[scr=rsfs,cal=euler]{mathalfa}
\usepackage{xfrac}
\usepackage[unicode,hidelinks,bookmarks=false]{hyperref}
\newcommand{\ie}{i.e.\ }
\newcommand{\cf}{cf.\ }
\newcommand{\resp}{respectively\ }
\newcommand{\Subsection}{\subsection}
\providecommand{\nobreakdash}{\nobreak\hbox{-}}
\setlength{\emergencystretch}{2em}
\multlinegap0pt
\usepackage{amssymb}
\newtheorem{theorem}{Theorem}
\newtheorem{lemma}{Lemma}
\newtheorem{proposition}{Proposition}
\title{The entropy of coclass one \\ and of SmallGroup(729,45)}
\author{Helga Boyer von Berghof}
\date{Source: arXiv:2609.06166v1 (CC BY 4.0)}
\begin{document}
\maketitle

\noindent\emph{Source and licence.} The entropy of coclass one \\ and of SmallGroup(729,45). Version arXiv:2609.06166v1 (CC BY 4.0), distributed by arXiv under the Creative Commons Attribution 4.0 International licence, http://creativecommons.org/licenses/by/4.0/, as recorded in the arXiv licence metadata for this exact version and shown on the abstract page https://arxiv.org/abs/2609.06166v1. Full licence notice: \texttt{licenses/arxiv-2609.06166-CC-BY-4.0.txt}.\par\smallskip

\noindent\textit{Editorial scope.} Counted unit: the article's theorem thm:CoClassOne (source0.tex lines 160--181). The article's complete original proof of it is delivered with it (source0.tex lines 182--293, one \texttt{proof} environment of 112 lines). No mathematics is written, repaired or re-derived: the statement and the proof are the article's own lines, in the article's own notation, and no retained line is substituted or re-flowed.  Retained uncounted as the source-local context the proof needs: lines 141--145; lines 300--309; lines 319--327.  Nothing was omitted from any retained span, so the package carries no omission record.  No retained line contains a citation, so the wrapper carries no bibliography and no bibliography entry.  No retained line contains a figure environment, an image-inclusion command or drawing code, so no image is delivered and the package declares no asset dependency.  Editorial formatting only, in addition: an AMS \texttt{amsart} preamble that restates the article's own declarations and macros used by the retained lines, namely the environments \texttt{theorem}, \texttt{lemma}, \texttt{proposition} on separate counters, together with the standard packages those lines need.  The retained environments print in this document as Theorem 1, Lemma 1, Proposition 1 in order of appearance, whereas the article prints the same items with the numbering of its own full document.  The complete native source of the article as distributed in the arXiv e-print for this exact version is bundled unchanged as \texttt{sources/209497/source0.tex}.
\begin{equation}
\label{eqn:RealMeasure}
{\bf P}_r(G)=\frac{y(G)^{g+1}}{\#\mathrm{Aut}(G)\cdot\#G}\cdot(p^g)^{g+1}\cdot\prod_{k=1}^{g}\,\left(1-\frac{1}{p^k}\right)^2
\cdot\left(1-\frac{1}{p^{g+1}}\right)\cdot\prod_{k=1}^{g+1-r}\,\left(1-\frac{1}{p^k}\right)^{-1}
\end{equation}

\begin{lemma}
\label{lem:FinitePart}
For each real number \(q\in\mathbb{R}\), \(q\ne 1\),
and non-negative integers \(K\le N\), the following formula holds
for finite partial sums of a geometric series:
\begin{equation}
\label{eqn:FinitePart}
q^K+q^{K+1}+\cdots+q^{N-1}+q^N=\frac{q^K-q^{N+1}}{1-q}.
\end{equation}
\end{lemma}

\begin{proposition}
\label{prp:GeometricVariant}
For each bounded real number \(q\in\mathbb{R}\), \(\lvert q\rvert<1\),
the sum of the following variant of the infinite geometric series is given by:
\begin{equation}
\label{eqn:GeometricVariant}
\sum_{n=1}^\infty\,n\cdot q^n=\frac{q}{(1-q)^2}.
\end{equation}
\end{proposition}

\begin{theorem}
\label{thm:CoClassOne}
The \textbf{entropy} of the descendant tree
\(\mathcal{T}(\Delta)\)
of all metabelian \(3\)-groups \(G\) of maximal nilpotency-class and fixed \textbf{coclass one},
\(\mathrm{cc}(G)=1\),
which are descendants of the extra-special \(3\)-group
\(\Delta=\langle 27,3\rangle\)
of exponent \(3\) without exceptions,
with respect to the normalized relative-measure
\({\bf P}_{rel}(G)=\frac{3^7}{2^7\cdot 13}\cdot{\bf P}_{r}(G)\)
associated to the real probability-measure in Formula
\eqref{eqn:RealMeasure},
is given by
\begin{equation}
\label{eqn:Entropy1}
H({\bf P}_{rel})
=-\sum\,{\bf P}_{rel}(G_n^i)\log({\bf P}_{rel}(G_n^i))
=\frac{1}{26}\biggl(41\log(3)-8\log(2)\biggr)
\approx{\bf 1.51915}.
\end{equation}
\end{theorem}

\begin{proof}
Using the Formula
\eqref{eqn:RealMeasure},
for \(r=g+1\),
and the information in Table
\ref{tbl:FormationLawsOne},
we calculate the measures
\({\bf P}_r(G)=\frac{y(G)^{2+1}}{\#\mathrm{Aut}(G)\cdot\#G}\cdot\frac{2^9\cdot 13}{3^3}\), grouped by TKTs: \\
a.1: \(\frac{(3^{n+1})^3}{2\cdot 3^{4n+4}\cdot 3^{2n+4}}\cdot\frac{2^9\cdot 13}{3^3}=\frac{2^8\cdot 13}{3^{3n+8}}\),
a.2: \(\frac{(3^{n+1})^3}{2\cdot 3^{4n+5}\cdot 3^{2n+4}}\cdot\frac{2^9\cdot 13}{3^3}=\frac{2^8\cdot 13}{3^{3n+9}}\),
a.3: \(\frac{(3^{n+1})^3}{2^2\cdot 3^{4n+4}\cdot 3^{2n+4}}\cdot\frac{2^9\cdot 13}{3^3}=\frac{2^7\cdot 13}{3^{3n+8}}\).

%\newpage
%--------------------------------------------------------------------------------
\renewcommand{\arraystretch}{1.2}
\begin{table}[ht]
\caption{Parametrized formation-laws of invariants, grouped by TKT}
\label{tbl:FormationLawsOne}
\begin{center}
\begin{tabular}{|c||c|c|c|}
\hline
 Transfer Kernel Type      & a.1, \(\varkappa\sim (0000)\)    &  a.2, \(\varkappa\sim (1000)\)   &  a.3, \(\varkappa\sim (2000)\)   \\
\hline
 State                     & \(n\ge 1\), \(3\) groups         & \(n\ge 0\), \(1\) group          & \(n\ge 0\), \(2\) groups         \\
 \(y(G_n^i)\)              & \(3^{n+1}\)                      & \(3^{n+1}\)                      & \(3^{n+1}\)                      \\
 \(\#\mathrm{Aut}(G_n^i)\) & \(2\cdot 3^{4n+4}\)              & \(2\cdot 3^{4n+5}\)              & \(2^2\cdot 3^{4n+4}\)            \\
 \(\#G_n^i\)               & \(3^{2n+4}\)                     & \(3^{2n+4}\)                     & \(3^{2n+4}\)                     \\
 \({\bf P}_r(G_n^i)\)      & \(\frac{2^8\cdot 13}{3^{3n+8}}\) & \(\frac{2^8\cdot 13}{3^{3n+9}}\) & \(\frac{2^7\cdot 13}{3^{3n+8}}\) \\
\hline
\end{tabular}
\end{center}
\end{table}
%--------------------------------------------------------------------------------

\noindent
Now we add all probability-measures
of metabelian \(3\)-groups of maximal nilpotency-class,
grouped by their TKTs,
and using the geometric series
\(\sum_{n=0}^{\infty}\,\frac{1}{27^n}
=\frac{1}{1-\frac{1}{27}}
=\frac{1}{\frac{26}{27}}
=\frac{27}{26}
=\frac{3^3}{2\cdot 13}\).
\begin{equation}
\label{eqn:a1}
3 \text{ times a.1}: \quad
\sum_{n=1}^{\infty}\,{\bf P}_r(G_n^1)=
\sum_{n=1}^{\infty}\,\frac{2^8\cdot 13}{3^{3n+8}}=
\frac{2^8\cdot 13}{3^{8}}\cdot\sum_{n=1}^{\infty}\,\left(\frac{1}{3^3}\right)^n=
\frac{2^8\cdot 13}{3^{8}}\cdot\frac{1}{2\cdot 13}=
\frac{2^7}{3^8},
\end{equation}
\begin{equation}
\label{eqn:a2}
1 \text{ times a.2}: \quad
\sum_{n=0}^{\infty}\,{\bf P}_r(G_n^2)=
\sum_{n=0}^{\infty}\,\frac{2^8\cdot 13}{3^{3n+9}}=
\frac{2^8\cdot 13}{3^{9}}\cdot\sum_{n=0}^{\infty}\,\left(\frac{1}{3^3}\right)^n=
\frac{2^8\cdot 13}{3^{9}}\cdot\frac{3^3}{2\cdot 13}=
\frac{2^7}{3^6},
\end{equation}
\begin{equation}
\label{eqn:a3}
2 \text{ times a.3}: \quad
\sum_{n=0}^{\infty}\,{\bf P}_r(G_n^3)=
\sum_{n=0}^{\infty}\,\frac{2^7\cdot 13}{3^{3n+8}}=
\frac{2^7\cdot 13}{3^{8}}\cdot\sum_{n=0}^{\infty}\,\left(\frac{1}{3^3}\right)^n=
\frac{2^7\cdot 13}{3^{8}}\cdot\frac{3^3}{2\cdot 13}=
\frac{2^6}{3^5}.
\end{equation}
Eventually, we add these three contributions,
taking into account their multiplicities, obtaining the 
\textbf{total measure of coclass one:}
\begin{equation}
\label{eqn:cc1}
\sum\,{\bf P}_{r}(G_n^i)=
3\cdot\frac{2^7}{3^8}+\frac{2^7}{3^6}+2\cdot\frac{2^6}{3^5}=
\frac{2^7}{3^7}\cdot(1+3+3^2)=
\frac{2^7\cdot 13}{3^7}=
\frac{\bf 1664}{\bf 2187}.
\end{equation}
With this result, we must normalize the probability-measure onto sum one, \({\bf P}_{rel}(G)=\frac{3^7}{2^7\cdot 13}\cdot{\bf P}_{r}(G)\):

%\newpage
%--------------------------------------------------------------------------------
\renewcommand{\arraystretch}{1.2}
\begin{table}[ht]
\caption{Normalized relative-measure and logarithms, grouped by TKT}
\label{tbl:RelativeMeasuresOne}
\begin{center}
\begin{tabular}{|c||c|c|c|}
\hline
 Transfer Kernel Type            & a.1, \(\varkappa\sim (0000)\)    &  a.2, \(\varkappa\sim (1000)\)   &  a.3, \(\varkappa\sim (2000)\)   \\
\hline
 State                           & \(n\ge 1\), \(3\) groups         & \(n\ge 0\), \(1\) group          & \(n\ge 0\), \(2\) groups         \\
 \({\bf P}_{rel}(G_n^i)\)        & \(\frac{2}{3^{3n+1}}\)           & \(\frac{2}{3^{3n+2}}\)           & \(\frac{1}{3^{3n+1}}\)           \\
 \(-\log({\bf P}_{rel}(G_n^i))\) & \((3n+1)\log(3)-\log(2)\)        & \((3n+2)\log(3)-\log(2)\)        & \((3n+1)\log(3)\)                \\
\hline
\end{tabular}
\end{center}
\end{table}
%--------------------------------------------------------------------------------

\noindent
We continue this proof of Theorem
\ref{thm:CoClassOne}
after an auxiliary Lemma
\ref{lem:FinitePart}
and Proposition
\ref{prp:GeometricVariant}.
\end{proof}

\end{document}
